


\subsection{The second watch - rough overview look at the video.}

\paragraph{We are looking for the following key things and moments:}\begin{enumerate}
    \item  SP1 (Stable point 1). Meaning the point where the pendulum is naturally resting and when there is no motion or external force working on it.
    \begin{enumerate}
        \item Pendulum in downward position.
        \item Natural swings about SP1.
    \end{enumerate}
    \item SP2 (stable point 2). When the pendulum is in upright position (180 degrees from SP1), this is a forced stabilized position.
    \begin{enumerate}
        \item Pendulum in upward position.
        \item Controller stabilizing.
    \end{enumerate}
    \item Disturbances.
    \begin{enumerate}
        \item When the pendulum is in SP2 and is disturbed externally.
    \end{enumerate}
\end{enumerate}


In the second watch we are to make timestamps for the intervals including all the 3 occurrences, they may appear more than once, as hinted to in the manual.
\begin{table}
    \centering
    \begin{tabular}{ccccc}
         Part &Short description.&  Notes&  Start frame (hr:min:sec:ms)& End frame (hr:min:sec:ms)\\
         SP1 &System introduction.&  First introduction to the pendulum system. The pendulum is resting at SP1&  00:00:08.008& 00:00:15.176\\
         SP1 &SP1 and swing.&  Second interval of the pendulum underneath horisontal-axis. The pendulum recieves a little push making the pendulum swing about SP1. The interval ends where as Kaya is moving the pendulum upwards and to the next stage of the system.&  00:00:19.887& 00:00:27.294\\
         SP2 (SP1 swings)&SP2 and release. Ends in natural swings&  Kaya continues to move the pendulum up towards SP2 and after letting go of the pendulum. The pendulums stored energy is decaying while osciliating, but then forced to stop by Kayas hand. Clip ends abruptly.&  00:00:27.327& 00:00:41.909\\
         SP2 &SP1 to SP2, controller on.&  The pendulum starts in SP1 while the controller is switched on. The student pushes the pendulum in order to make the controller act. The controller progressively increases the pendulums swing until it reaches the other side of the horisontal-axis. The swing becomes increasingly calm when reaching SP2&  00:00:41.942& 00:01:09.937\\
         Disturbances 1 &SP2 controller drop.&  First disturbance. Looks like the student is turning off the controller. The pendulum is falling, decaying and then the student catches it making the pendulum fully stop.&  00:01:13.974& 00:01:27.254\\
         SP2 &SP1 to SP2, controller on and with issues.&  The student releases the pendulum several degrees away from SP1 and turns on the controller. The controller is having issues caused by the lack of room and the amount of noise in the swing. This interval ends when the pendulum has reached SP2 and is almost not moving.&  00:01:29.823& 00:01:54.581\\
         Disturbances 2-5 &Light taps on the pendulum.&  The student lightly taps the side of the pendulum multiple times. We can clearly see the controller compensating for the force being applied for each event. The video ends here.&  00:01:54.620& 00:02:03.123\\
          &&  &  & \\
    \end{tabular}
    \caption{Observations after first watch}
    \label{tab:placeholder}
\end{table}
\subsubsection{Notes after second watch:}
Some of the segments might not be suitable for taking measurements like the first SP1 segment. Do we need all 4 last disturbance events, are they alike?



\section{Taking measurements}
Description, timestamp, name or something to indicate for which segment the measurements are taken from.
\begin{table}
    \centering
    \begin{tabular}{ccc}
           Variable name/ Description& Symbol &Value\\
           Release time&  $t_0$&\\
           Initial angular displacement&   $\tilde\theta_0$&\\
           Direction of initial motion&  -&\\
           Peak angle 1 (first swing)&  $\tilde\theta_1$&\\
           Peak angle 2 (second swing)&  $\tilde\theta_2$&\\
           Peak angle 3 (third swing)&  $\tilde\theta_3$&\\
 Time for peak 1& $t_1$&\\
 Time for peak 2& $t_2$&\\
 Time for peak 3& $t_3$&\\
 Number of complete oscillations& -&\\
 Visible cart motion, if any& -&\\
 Settling time, if visible& $t_s$&\\
 Final pendulum position& -&\\
 Downward equilibrium& $\tilde\theta$&0\\
    \end{tabular}
    \caption{Natural motion SP1 event 1}
    \label{tab:placeholder}
\end{table}
Description, timestamp, name or something to indicate for which segment the measurements are taken from.
\begin{table}
    \centering
    \begin{tabular}{ccc}
           Variable name/ Description& Symbol &Value\\
           Release time&  $t_0$&\\
           Initial angular displacement&   $\tilde\theta_0$&\\
           Direction of initial motion&  -&\\
           Peak angle 1 (first swing)&  $\tilde\theta_1$&\\
           Peak angle 2 (second swing)&  $\tilde\theta_2$&\\
           Peak angle 3 (third swing)&  $\tilde\theta_3$&\\
 Time for peak 1& $t_1$&\\
 Time for peak 2& $t_2$&\\
 Time for peak 3& $t_3$&\\
 Number of complete oscillations& -&\\
 Visible cart motion, if any& -&\\
 Settling time, if visible& $t_s$&\\
 Final pendulum position& -&\\
 Downward equilibrium& $\tilde\theta$&0\\
    \end{tabular}
    \caption{Natural motion SP1}
    \label{tab:placeholder}
\end{table}

\textbf{Description, timestamp, name or something to indicate for which segment the measurements are taken from.}
\begin{table}
    \centering
    \begin{tabular}{ccc}
         Variable name/ Description&  Symbol & Value\\
         Beginning of controlled operation&  $t_on$& \\
         Maximum positive angular deviation&  $\tilde\theta_{max}$& \\
         Maximum negative angular deviation&  $\tilde\theta_{min}$& \\
         Maximum cart position&  $x_{max}$& \\
         Minimum cart position&  $x_{min}$& \\
         Stabilization time&  $t_{stab}$& \\
         Direction of first corrective cart movement&  -& \\
         Remaining vibration around upright point&  -& \\
 Upright position& $\hat\theta$&0\\
    \end{tabular}
    \caption{Controlled operation around SP2 - Take 1}
    \label{tab:placeholder}
\end{table}
Description, timestamp, name or something to indicate for which segment the measurements are taken from.
\begin{table}
    \centering
    \begin{tabular}{ccc}
         Variable name/ Description&  Symbol & Value\\
         Disturbance time&  $t_d$& \\
         Maximum angular deviation after disturbance&  $\tilde\theta_{max}$& \\
         Maximum cart displacement&  $\Delta x_{max}$& \\
         Recovery time&  $t_r$& \\
         Whether cart correction remains visible&  -& \\
 Whether pendulum leaves $\pm 20^\circ$ region& -&\\
 Whether pendulum returns to SP2& -&\\
 Final state& -&\\
    \end{tabular}
    \caption{Disturbance event 1}
    \label{tab:placeholder}
\end{table}
Description, timestamp, name or something to indicate for which segment the measurements are taken from.
\begin{table}
    \centering
    \begin{tabular}{ccc}
         Variable name/ Description&  Symbol & Value\\
         Disturbance time&  $t_d$& \\
         Maximum angular deviation after disturbance&  $\tilde\theta_{max}$& \\
         Maximum cart displacement&  $\Delta x_{max}$& \\
         Recovery time&  $t_r$& \\
         Whether cart correction remains visible&  -& \\
 Whether pendulum leaves $\pm 20^\circ$ region& -&\\
 Whether pendulum returns to SP2& -&\\
 Final state& -&\\
    \end{tabular}
    \caption{Disturbance event 2}
    \label{tab:placeholder}
\end{table}
Description, timestamp, name or something to indicate for which segment the measurements are taken from.
\begin{table}
    \centering
    \begin{tabular}{ccc}
         Variable name/ Description&  Symbol & Value\\
         Disturbance time&  $t_d$& \\
         Maximum angular deviation after disturbance&  $\tilde\theta_{max}$& \\
         Maximum cart displacement&  $\Delta x_{max}$& \\
         Recovery time&  $t_r$& \\
         Whether cart correction remains visible&  -& \\
 Whether pendulum leaves $\pm 20^\circ$ region& -&\\
 Whether pendulum returns to SP2& -&\\
 Final state& -&\\
    \end{tabular}
    \caption{Disturbance event 3}
    \label{tab:placeholder}
\end{table}
\section{Calculations}